\section{Binary interval orbits and supplied normal-surface topology}
\label{sec:native-interval-orbits}

The quotient-kernels delivery in \texttt{6cfd26c80} supplies an implementation
of the classical Agol--Hass--Thurston orbit algorithm, with a sharper periodic
merger. We integrate its interval kernel, normal-coordinate adapter, fixtures
and tests. This closes the unweighted orbit-counting implementation gap
identified in the earlier normal-surface reviews. It does not close the
surface-search, weighted component-extraction, cutting, attachment or
diagram-to-exterior provenance gaps. Ordinary recognition does not yet
dispatch this supplied-surface interface.

The integration provenance records the archive's full SHA-256.
All 84 manifest entries verify.
Seven integrated source/test files match
the reviewed delivery exactly, as recorded in
\path{data/quotient-orbit-integration.json}. The coefficient-span and joint
cyclic-overlap changes in the same delivery are separate, pending work.
This section reports our reproduction, not the delivery's historical
full-suite timings or its measurements against an older checkout.

\subsection{The encoded relation and a sharp local merger}

An interval pairing is an isometry from $[a,b]$ to $[c,d]$, with inclusive
integer endpoints, equal positive widths, and $a\le c$. Its map is either
$x\mapsto x+c-a$ or $x\mapsto a+d-x$. Endpoints have $O(\log(N+1))$ bits;
the interface stores $k$ such records on $[0,N-1]$, never $N$ vertices.
The requested answer is the number of equivalence classes generated by
these identifications, including uncovered singleton points.

A positive translation with $a<c\le b+1$ has period $p=c-a$ and support
$[a,d]$ of size at least $2p$. Its classes are precisely the residues
modulo $p$ within this support. Consider two such supports, of periods
$p,q$, with $h$ points of intersection, and put $g=\gcd(p,q)$.

\begin{proposition}
The relation generated by the two pairings equals the residue relation
modulo $g$ on their union if and only if $h\ge p+q-g$.
\end{proposition}
\begin{proof}
Every identification preserves residue modulo $g$. Color each point of the
intersection by its generated orbit. This sequence has periods $p$ and $q$.
The Fine--Wilf theorem gives period $g$ at length $p+q-g$
\cite{fine-wilf-orbits}. This length is at least $\max(p,q)$, so every residue
class in either support meets the intersection. Thus every point reaches
the appropriate common $g$-class, proving sufficiency.

For necessity, first replace each support separately by its residue
classes: there are $p+q$ vertices. Each shared point supplies one edge
between these two sets of vertices. These $h$ edges can reduce the number
of components by at most $h$, leaving at least $p+q-h$. Below the proposed
threshold this exceeds $g$. Both supports contain every residue modulo
$g$, so equality with the $g$-residue relation is impossible.
\end{proof}

The condition is exact for these two generators alone. Additional pairings
can identify further classes; failure of this test is not a disconnectedness
certificate. A justified replacement remains valid when any common set of
other generators is adjoined. The maintained default uses this threshold;
\texttt{periodic\_rule='aht'} retains $p+q$ as a comparison mode.

\subsection{Counting without expanding integer points}

\path{../fast/fastunknot/interval_orbits.py} repeatedly deletes identities,
contracts uncovered intervals while adding their singleton counts, trims
duplicate reflection arrows, merges qualifying periodic supports, transmits
pairings through a maximal rightmost carrier, and truncates its exclusive
suffix. Trimming a reflection leaves its fixed midpoint unpaired; the next
static contraction counts it if no other generator meets it. Truncation
removes points already identified with surviving representatives and adds
no singleton count. Transmission computes maximal translation powers by
integer division, rather than executing one shift per point. These are the
classical AHT operations \cite{aht-orbits}; their order and maximality
conditions are part of the implementation contract.

For the sharper merger, let $w_1,w_2$ be the old domain widths. The merged
width is
\[
 w'=w_1+w_2+p+q-h-g\le w_1+w_2\le 2w_1w_2.
\]
Consequently the AHT potential $4^k\prod_i w_i$ decreases by at least a
factor two at each such merge. When no sharp merge is available, no
classical merge is available either. The remaining maximal-transmission
progress argument is therefore unchanged. There are polynomially many
cycles in $k$ and $\log(N+1)$; sorting, scanning all candidate pairs and
binary integer arithmetic within each cycle preserve a polynomial bit
bound. We do not infer an optimal polynomial degree from these simple
list-based data structures.

For $r$ adjacent period-$p$ pairings on $(r+1)p$ points,
\[
 [ip,(i+1)p-1]\longrightarrow[(i+1)p,(i+2)p-1],\qquad 0\le i<r,
\]
the count is $p$. The classical threshold misses every neighboring overlap
of size $p$, and the rightmost suffix is removed one block per cycle: there
are $r+1$ cycles including the final static contraction. The sharp rule
merges the entire chain in one cycle and finishes in two. Candidate-pair
scanning still has a cost; cycle savings alone are not timing evidence.

\texttt{max\_cycles} limits begun cycles. It is neither a bit-operation
allowance nor a wall-clock deadline. Exhaustion returns an incomplete
result with no count. Cooperative callbacks propagate cancellation. The
signed wrapper constructs $2k$ pairings on a two-sheeted interval of size
$2N$, with an independent parity bit on each identification. If its orbit
count is $C_2$ and the base count is $C$, then $C_2-C$ components have a
consistent parity labeling and $2C-C_2$ do not. The two counts share one
cycle allowance. Order reversal of an interval does not by itself specify
this parity character. A same-orbit query compares counts before and after
adjoining a singleton identification; it does not extract a component.

\subsection{From normal coordinates to three orbit queries}

\path{../fast/fastunknot/normal_surface_orbits.py} accepts a finite connected
compact orientable tetrahedral manifold with exactly one torus boundary,
and seven nonnegative binary coordinates per tetrahedron. It reconstructs
reciprocal face identifications, edge orientations and links, vertex links,
ambient orientation and boundary topology. It rejects reversed edge
self-identifications, invalid links and unsupported ambient manifolds.
Coordinates must satisfy matching equations and the quadrilateral
constraint. They need not be primitive, fundamental or extremal.

Allocate one interval for the intersection points on each global edge.
The ordered points nearest a face corner are determined by its triangle
coordinate plus the appropriate quadrilateral coordinate. Each face has
at most three such arc families, represented by interval isometries.
Local edge directions are transported through the actual face gluings;
global vertex labels cannot determine direction in a one-vertex manifold.
One copy of each paired face is used. This gives $O(t)$ pairings, with
at most $4W$ points when the surface has $W$ normal discs. The normal
surface has the same components as this graph: every attached disc has a
connected boundary already in the graph. A second graph uses only boundary
edges and faces, avoiding spurious isolated interior points.

The three queries count the surface, the surface with all coordinates
doubled, and its boundary. In an orientable ambient manifold the doubled
surface is the horizontal boundary of its normal interval bundle, hence
its orientation double. If the first two counts are $C,C_2$, the numbers
of orientable and nonorientable components are $C_2-C$ and $2C-C_2$.
This geometric doubling handles one-sided components correctly: two copies
of the coordinate vector of a M\"obius band describe an annulus.
Total Euler characteristic is computed directly as $V-E+F$ from global
edge intersections, normal arcs and discs. For connected output, boundary
count and Euler characteristic determine genus or crosscap number.
For disconnected output these are aggregate statistics; no distribution
of genera among components is asserted.

A connected surface with $\chi=1$ and one boundary curve is a disc.
To test whether that disc compresses the torus, the adapter checks the
boundary intersection cocycle over $\mathbb F_2$. Solving for vertex
potentials tests whether this cocycle is a coboundary. An inconsistency
means nonzero boundary homology. A simple closed curve on a torus is
essential exactly when this mod-two class is nonzero: an essential curve
has a primitive integral slope. This last assertion uses simplicity and
the torus hypothesis, and is not a general multicurve criterion.

All three orbit calls share one cycle allowance. If any stops early, the
adapter returns \texttt{INCONCLUSIVE} without a disc assertion. Completed
topology has polynomial cost in $t$ and $\log(W+1)$. A returned compressing
disc assertion concerns the supplied manifold only. It does not establish
that this manifold is the exterior of the user's knot, and it is not a
new independent serialized knot certificate.

\subsection{Validation and controlled measurements}

The 22 focused tests pass with Regina. They include exhaustive small
pairing systems, random literal union-find comparisons, both merger
thresholds, binary endpoint inflation, signed parity graphs, cancellation
and shared budgets. The geometric tests compare layered meridians with
literal arc graphs, compare admissible normal-surface sums and component
types with Regina, exhaust one-tetrahedron one-face-pairing manifold
validation, and randomize tetrahedron and vertex labels. They also check
huge parallel discs, odd and even multiples of a one-sided surface,
boundary-parallel discs, and a 128-tetrahedron meridian with more than
$10^{26}$ normal discs. Detailed logs are retained under
\path{data/quotient-orbit-*}. The complete maintained suite passes all 805
tests with Regina in 191.413 seconds.

\begin{center}\small
\begin{tabular}{@{}lrrrrr@{}}
\toprule
Query & Classical & Sharp & A/A & Ratio & Cycles \\
\midrule
Chain, 8 pairings & 0.414 & 0.085 & 0.086 & 4.871 & 9/2 \\
Chain, 16 pairings & 2.212 & 0.162 & 0.162 & 13.066 & 17/2 \\
Chain, 32 pairings & 12.904 & 0.317 & 0.315 & 41.464 & 33/2 \\
Chain, 64 pairings & 89.518 & 0.602 & 0.609 & 147.258 & 65/2 \\
Chain, 128 pairings & 759.598 & 1.195 & 1.199 & 649.887 & 129/2 \\
Meridian, 4 tetrahedra & 1.907 & 1.971 & 1.954 & 0.971 & 23/25 \\
Meridian, 32 tetrahedra & 73.945 & 70.930 & 71.129 & 1.038 & 107/109 \\
Meridian, 128 tetrahedra & 2854.560 & 2862.672 & 2822.536 & 0.982 & 395/397 \\
$2^{2000}$ parallel discs & 1.956 & 2.059 & 2.061 & 0.952 & 20/22 \\
\bottomrule
\end{tabular}
\end{center}
Median milliseconds; ratio is the median paired classical/sharp time. Cycles are classical/sharp. Surface queries include validation and all three orbit counts; chain queries count interval orbits only.


The maintained driver \path{../fast/benchmark_orbits.py} uses seed 261008171,
one warmup and five measured rounds with shuffled classical, sharp and
identical-sharp control arms. There are 29 interval cases and 13 surface
cases, giving 126 warmup calls and 630 measured calls. Source hashes agree
before and after the run, pinned to the integrated code at \texttt{ac0dfceb9}.
Raw samples, arm orders, operation counts and complete last outputs are in
\path{../fast/results/orbits_20261008.json}.

At 128 adjacent pairings the paired classical/sharp ratio is 649.887,
with 129 versus two cycles. This is a constructed interval relation,
not a knot-diagram workload. Six independently generated tiny parity
graphs also determine expected counts after replacing each vertex by
$2^b+1$ ordered points, for $b=16,256,4096,16384$. A component with
consistent reversal parity lifts to $2^b+1$ components; an inconsistent
one lifts to $2^{b-1}+1$. All expanded counts agree without materializing
the large graph. Five templates have essentially unchanged cycle counts;
one drops from nine to seven, with measured ratios from 1.126 to 1.211.

For actual layered-torus meridians with 4--128 tetrahedra, complete topology
takes about 2 ms to 2.9 seconds. The sharper rule has paired ratios
0.971--1.038, and takes two more total cycles on these examples. Earlier
valid merges can change later scheduling, so a sharper local threshold
does not monotonically reduce the total cycle count. At 128 tetrahedra
both modes validate the surface with more than $10^{26}$ discs. On three
tetrahedra with $2^{2000}$ parallel meridians, the sharp query takes about
2.06 ms and 22 cycles. Across these parallel-copy cases it is roughly
4--7\% slower than the classical mode. Identical-control paired ratios
range from 0.982 to 1.030 across all cases.

Explicit connectivity checks run only below a declared 100,000-point cap;
larger cases are omitted, not described as timeouts. They perform less work
than the encoded surface query, which includes manifold validation and
three counts, so their single samples are diagnostic rather than a matched
speedup comparison. The evidence establishes useful encoded local queries
and a sharp-merger gain on its target family. It establishes neither a
uniform local speedup nor a whole-recognizer speedup.

\subsection{Repairing worker communication}

The integration also repairs Regina worker communication. A retry with
\texttt{communicate(None)} after timing out on a large request can
leave input unsent on supported Python implementations. The worker then
waits for data while its parent waits for completion. The revised code
stages UTF-8 input, output and diagnostics in temporary files, waits
cooperatively, and reads the completed streams. An exit stack closes every
stream; cancellation kills and reaps a live child. Startup, staging, waiting
and reading share the allowance through checks before and after blocking
operations. Filesystem calls are not made preemptible by these checks.

All nine focused worker tests pass, including delayed reading of a large
request, simultaneous large UTF-8 output and diagnostics, and local/global
cancellation cleanup. This is a reliability repair to an existing complete
backend, separate from the new supplied-surface API.

\subsection{Christoffel delivery: initial review before the native contract}

The adjacent Christoffel archive has SHA-256 prefix
\texttt{9dc35414f591da0b}; its full hash and reproduction status are recorded
in \path{data/quotient-orbit-review.json}.
All 32 manifest entries verify; its 22 tests and both saved example
replays pass locally. The report supplies a compressed primitive-power
test for freely and cyclically reduced rank-two words, with independent
local replay. For coherently signed exponent counts $(a,b)$, divide by
$d=\gcd(|a|,|b|)$ to obtain primitive counts $(u,v)$ after sign
normalization. Prefix weights $(v,-u)$ characterize Christoffel powers
by width $u+v-1$, including the pure-letter case. Concatenation summaries
store total displacement and minimum/maximum prefix displacement, so
the exponentially long word need not be materialized.

The useful group-theoretic consequence is conditional on provenance:
in a torsion-free group a relation $P^d=1$ with $P$ primitive kills $P$.
For a two-generator knot group this leaves a cyclic group, necessarily
$\mathbb Z$ by knot-group abelianization. The resulting unknot conclusion
requires a verified presentation trace from the input knot. The delivery's
small braid demonstration has an independent literal, capped trace replay;
its higher-rank projection prototype lacks an independent whole-round
checker and assumes torsion-freeness. Its WordArena adapter is tested
against a schema stub, not our maintained search.

At this initial checkpoint the compressed verifier required its existing
rank-one/empty-relator terminal. The later integration in
Section~\ref{sec:primitive-power-terminal} adds a separate version-five
terminal and independent source-bound replay, while preserving the old
contract. The report's
conditional depth estimate, polynomial initial cost times $2^{O(d)}$,
still needs a controlled presentation producer, a depth bound and a
guarantee of reaching a decisive endpoint. Neither this estimate nor the
polynomial orbit kernel proves unrestricted quasi-polynomial recognition.

During this review, the original deliveries were archived unchanged as
reports 45 (Christoffel) and 46 (quotient kernels); the incoming ZIP paths
in the provenance record refer to their arrival commit. The newer report
47 supplies a separate orbit implementation with independent local-rule
replay, marked-component incidence and a two-meridian search proposal.
Its 512-entry manifest verifies. Reconciliation of that certificate API
with the sharper maintained kernel is subsequent work, not validation
claimed for this checkpoint.
Section~\ref{sec:orbit-certificates} subsequently integrates independent
orbit replay and marked-union incidence, while retaining the geometric
and recognition boundaries above.
